\chapter{The Risk Assessment Process — An Overview}
\label{chap:The Risk Assessment Process — An Overview}

%---------------------------- SECTION ----------------------------
\section{From Theory to Practice}
\paragraph{The first six chapters of this book have covered a considerable amount of ground. We have defined risk, explored what risk assessment is and why it matters, traced its history, mapped the risk landscape, examined the frameworks that structure it, and surveyed the regulatory obligations that make it a professional imperative. That foundation is important — but it is, ultimately, preparation for the work itself.}
\paragraph{This chapter marks the transition from understanding to doing.}
\paragraph{Here, we bring together everything covered so far and lay out the risk assessment process as a coherent, end-to-end sequence of activities. Think of it as the map of the journey ahead — a high-level view of the full process before we descend into the detail of each step in Chapters~\ref{chap:Step 1 — Identifying Hazards and Risk Sources} through~\ref{chap:Step 5 - Recording, Reviewing, and Monitoring}.}
\paragraph{It also addresses something that the preceding chapters have touched on but not yet examined directly: how does the risk assessment process fit within an organisation's broader governance and compliance arrangements? Who owns it? How is it structured? How does it connect to the decisions that matter? These are questions that compliance and legal professionals are frequently called upon to answer — and getting the answers right makes the difference between a risk assessment process that genuinely serves the organisation and one that exists primarily on paper.}

%---------------------------- SECTION ----------------------------
\section{The Five Steps — A Framework for the Journey}
\paragraph{The risk assessment process, at its core, follows a logical and repeatable sequence. Whilst different frameworks and methodologies use slightly different language, the underlying logic is broadly consistent — and maps closely to the process outlined in ISO 31000, which we examined in Chapter~\ref{chap:Risk Management Frameworks}.}
\paragraph{For the purposes of this book, we will organise the process around \textbf{five key steps}:}
\begin{enumerate}
    \bitem{Identifying hazards and risk sources}{What could go wrong? Where do the risks come from?.}
    \bitem{Determining who or what is affected}{Who or what is exposed to those risks, and in what ways?.}
    \bitem{Evaluating and analysing risks}{How likely are these risks to materialise, and how serious would the consequences be?.}
    \bitem{Deciding on controls and actions}{What should we do about the risks we have identified and evaluated?.}
    \bitem{Recording, reviewing, and monitoring}{How do we document what we have done, and how do we ensure the assessment remains current and effective?.}
\end{enumerate}
\paragraph{These five steps are not a rigid checklist to be worked through mechanically and then set aside. They are a cycle — each step informing the next, and the whole process returning to the beginning as circumstances change, new risks emerge, and existing controls are tested by experience.}
\paragraph{Before examining the steps themselves, it is worth pausing on two things: the activities that should happen \textit{before} the five steps begin, and the activities that run \textit{alongside} the process throughout.}

%---------------------------- SECTION ----------------------------
\section{Before You Start — Scope, Context, and Criteria}
\paragraph{One of the most common reasons risk assessments fail to deliver their potential is that insufficient thought is given to the work that needs to happen before the formal process begins. Jumping straight into hazard identification without first establishing clear parameters is a reliable way to produce an assessment that is either too broad to be useful, too narrow to be complete, or simply aimed at the wrong things.}
\paragraph{The preparatory work falls into three areas...}

\subsection{Defining the Scope}
\paragraph{As we noted in Chapter~\ref{chap:What Is a Risk Assessment}, scope is one of the most important decisions in any risk assessment. At this stage, you need to be clear about:}
\begin{itemize}
    \bitem{What is being assessed}{a specific process, a business function, a product or service, a project, or the organisation as a whole.}
    \bitem{What is not being assessed}{explicitly acknowledging what falls outside the scope is as important as defining what falls within it, because it prevents false assurance and manages expectations.}
    \bitem{The time horizon}{is this assessment concerned with immediate risks, medium-term exposures, or longer-term strategic vulnerabilities?}
    \bitem{The organisational boundaries}{does the assessment cover only the organisation itself, or does it extend to third parties, supply chains, or other entities whose activities create exposure?}
\end{itemize}
\paragraph{Getting scope wrong — either too broad or too narrow — is one of the most common sources of risk assessment failure. A scope that is too broad produces an unwieldy assessment that cannot be maintained or acted upon. A scope that is too narrow creates blind spots that can prove costly when risks emerge from areas that were not considered.}

\subsection{Understanding the Context}
\paragraph{Every risk assessment takes place within a context — an internal and external environment that shapes what risks are relevant, how significant they are, and what responses are available. Understanding that context before beginning the assessment produces better results.}
\paragraph{\textbf{External context} includes the regulatory environment, the competitive landscape, the economic climate, the technological environment, and the broader societal and political factors that may affect the organisation's risk profile. For compliance professionals, the regulatory dimension of external context is particularly significant — changes in regulation, new enforcement priorities, and evolving supervisory expectations can fundamentally alter the risk landscape.}
\paragraph{\textbf{Internal context} includes the organisation's objectives, its governance structure, its culture, its existing risk management arrangements, its capabilities and resources, and its history of incidents and near-misses. Understanding what the organisation is trying to achieve — and the constraints within which it is operating — is essential to conducting a risk assessment that is genuinely relevant and actionable.}

\subsection{Establishing Risk Criteria}
\paragraph{\textbf{Risk criteria} are the standards against which the significance of risks will be evaluated. They answer the question: \textit{how do we decide whether a risk is acceptable, requires action, or demands urgent attention?}}
\paragraph{Risk criteria typically encompass:}
\begin{itemize}
    \bitem{Likelihood criteria}{what scale do we use to assess how probable a risk is? What does "likely" mean in our context — once a year, once a decade, a particular percentage probability?}
    \bitem{Consequence criteria}{what scale do we use to assess the severity of impact? What level of financial loss, regulatory penalty, reputational damage, or operational disruption constitutes a "minor" impact versus a "major" one?}
    \bitem{Risk appetite}{what overall level of risk is the organisation willing to accept in pursuit of its objectives? Which risk types are subject to zero tolerance, and which can be accepted within defined limits?}
\end{itemize}
\paragraph{Establishing clear, shared risk criteria before beginning the assessment prevents one of the most persistent problems in risk assessment practice: different people using different implicit scales to evaluate the same risks, producing results that are inconsistent and incomparable. We will explore risk criteria and risk appetite in more detail in Chapter~\ref{chap:Step 3 — Evaluating and Analysing Risks}.}

%---------------------------- SECTION ----------------------------
\section{Alongside — Communication, Consultation, and Governance}
\paragraph{Two activities run alongside the five-step process throughout — not as discrete steps themselves, but as continuous threads that weave through the entire assessment.}

\subsection{Communication and Consultation}
\paragraph{A risk assessment conducted in isolation — by a single individual or a small team working without input from the wider organisation — is almost always less complete and less credible than one that draws on diverse perspectives and genuine engagement.}
\paragraph{\textbf{Communication} ensures that the right people are informed about the assessment — its purpose, its scope, its progress, and its findings. This includes both the people contributing to the assessment and those who will need to act on its outputs.}
\paragraph{\textbf{Consultation} goes further — it actively seeks the views, knowledge, and experience of stakeholders who have relevant insight. This might include operational staff who understand the day-to-day realities of a process, subject matter experts who can assess technical risks, senior leaders who can provide strategic context, or external advisers who bring specialist knowledge.}
\paragraph{For compliance professionals, consultation is particularly important for two reasons. First, it produces better assessments — the people closest to a risk are often the best placed to identify it and understand its nature. Second, it builds ownership — people who have been genuinely involved in identifying and assessing risks are more likely to engage seriously with the controls and actions that follow.}

\subsection{Governance and Oversight}
\paragraph{The risk assessment process needs to be governed — there need to be clear answers to questions about who has authority to make risk-related decisions, who is responsible for overseeing the process, and how risk assessment outputs connect to the organisation's broader governance arrangements.}
\paragraph{Key governance questions include:}
\begin{itemize}
    \item{\textbf{Who owns the process?} Who is responsible for ensuring that risk assessments are conducted, documented, and reviewed?}
    \item{\textbf{Who has decision-making authority?} For risks above a certain threshold, who has the authority to decide on the appropriate response — and who needs to be consulted or informed?}
    \item{\textbf{How does the process connect to the board and senior management?} How are material risks escalated, and how does the board maintain adequate oversight without being overwhelmed with detail?}
    \item{\textbf{How is the process quality-assured?} Who provides independent challenge to the assessment — ensuring that it is not simply a reflection of what people want to hear?}
\end{itemize}
\paragraph{These governance questions are not abstract. They have direct implications for how the risk assessment process is designed, how it is resourced, and how credible it will appear to regulators and other external stakeholders. We will return to governance in depth throughout Part III.}

%---------------------------- SECTION ----------------------------
\section{Step One — Identifying Hazards and Risk Sources}
\paragraph{The first step is to identify, as comprehensively as possible, the hazards and risk sources relevant to the scope of the assessment. This means looking systematically at the activities, processes, relationships, and environments covered by the assessment and asking: \textit{what could go wrong here?}}
\paragraph{Effective hazard identification draws on a range of sources and techniques — structured workshops, process mapping, incident and near-miss data, regulatory guidance, industry benchmarks, and the knowledge of experienced staff. It requires both analytical rigour and creative thinking — the willingness to consider not just the risks that have materialised before, but those that could plausibly arise in the future..}
\paragraph{The output of this step is a comprehensive list of identified risks — the raw material on which the rest of the process operates. We will explore the techniques and practicalities of hazard identification in detail in Chapter~\ref{chap:Step 1 — Identifying Hazards and Risk Sources}.}

%---------------------------- SECTION ----------------------------
\section{Step Two — Determining Who or What Is Affected}
\paragraph{Having identified the hazards and risk sources, the next step is to think carefully about who or what is exposed to those risks — and in what ways. This step is sometimes treated as a minor adjunct to hazard identification, but it deserves more attention than that.}
\paragraph{Understanding exposure shapes the subsequent evaluation of risk. The same hazard can present very different levels of risk depending on who is exposed to it and in what circumstances. A data security vulnerability that affects a small number of internal users presents a very different risk profile to one that affects thousands of customers' personal financial data.}
\paragraph{For compliance and legal professionals, this step also has direct relevance to regulatory obligations — many regulatory requirements are specifically framed around the protection of particular groups of people or interests, and ensuring that those groups are explicitly considered in the risk assessment is both good practice and a direct component of regulatory compliance. Chapter~\ref{chap:Step 2 — Determining Who or What Is Affected} explores this step in full.}

%---------------------------- SECTION ----------------------------
\section{Step Three — Evaluating and Analysing Risks}
\paragraph{With a clear picture of what the risks are and who they affect, the third step is to evaluate their significance — to develop an informed view of how likely each risk is to materialise, how serious the consequences would be if it did, and how those factors combine to produce an overall assessment of each risk's priority.}
\paragraph{This is the analytical heart of the risk assessment process. It involves applying the risk criteria established at the outset — likelihood scales, consequence scales, and risk appetite — to each identified risk, and producing an evaluated risk profile that can inform decision-making and prioritisation.}
\paragraph{Evaluation can be qualitative — using descriptive judgements about likelihood and consequence — or quantitative, using numerical data and probabilistic modelling. In practice, most risk assessments in compliance and legal contexts use a combination of both, applying quantitative techniques where reliable data exists and qualitative judgement where it does not.}
\paragraph{The output of this step is typically a \textbf{risk register} — a structured record of identified and evaluated risks that forms the backbone of the organisation's risk management documentation. We will explore evaluation techniques and the construction of risk registers in detail in Chapter~\ref{chap:Step 3 — Evaluating and Analysing Risks}.}

%---------------------------- SECTION ----------------------------
\section{Step Four — Deciding on Controls and Actions}
\paragraph{Evaluation tells you where the risks are and how significant they are. The fourth step is to decide what to do about them.}
\paragraph{This involves selecting and implementing controls — measures designed to reduce the likelihood of a risk materialising, to limit its consequences if it does, or both. The range of possible responses is wide: risks can be avoided, reduced, transferred, or consciously accepted. The right response for any given risk depends on its nature, its significance, the options available, and the organisation's risk appetite.}
\paragraph{For compliance professionals, this step is where the risk assessment process connects most directly to the organisation's day-to-day operations. Controls are not abstract — they are specific policies, procedures, systems, and behaviours that need to be designed, implemented, maintained, and monitored. The quality of the controls decided upon in this step determines, more than anything else, whether the risk assessment actually delivers the protection and compliance outcomes it is designed to achieve.}
\paragraph{We will explore the full range of control options and how to choose between them in Chapter~\ref{chap:Step 4 — Deciding on Controls and Actions}.}

%---------------------------- SECTION ----------------------------
\section{Step Five — Recording, Reviewing, and Monitoring}
\paragraph{The final step — though it is perhaps more accurate to think of it as an ongoing activity than a discrete step — encompasses the documentation, review, and monitoring activities that ensure the risk assessment remains current, effective, and evidentially valuable.}
\paragraph{\textbf{Recording} means capturing the process and its outputs in a form that is clear, accessible, and defensible — creating the documentary evidence that demonstrates to regulators, auditors, and other stakeholders that the organisation has conducted a genuine and proportionate risk assessment.}
\paragraph{\textbf{Reviewing} means returning to the assessment periodically — and in response to specific triggers — to ensure that it remains accurate and relevant. Risks change. Controls fail. New hazards emerge. An assessment that is not reviewed becomes outdated, and an outdated assessment can provide false assurance that is, in some respects, more dangerous than no assessment at all.}
\paragraph{Monitoring means maintaining ongoing visibility of the risk landscape between formal review cycles — watching for early warning signs that risks are changing, that controls are not working as intended, or that new exposures are developing. We will cover all of this in Chapter~\ref{chap:Step 5 - Recording, Reviewing, and Monitoring}.}

%---------------------------- SECTION ----------------------------
\section{How the Process Fits Within the Organisation}
\paragraph{Understanding the five steps is important. Understanding how they fit within the organisation's broader governance and compliance arrangements is equally so — because a risk assessment process that operates in isolation from the rest of the organisation's decision-making will never deliver its full potential.}

\subsection{Connecting to Strategy and Objectives}
\paragraph{Risk assessment should be directly connected to the organisation's strategy and objectives. Risks are only meaningful in the context of what the organisation is trying to achieve — and a risk assessment that is not anchored to clear organisational objectives risks becoming an abstract exercise that produces lists of risks without generating genuine insight or action.}
\paragraph{For compliance professionals, this connection is particularly important in conversations with boards and senior management. Framing risk assessment in terms of its relevance to strategic objectives — rather than purely as a regulatory obligation — is one of the most effective ways to secure genuine leadership engagement with the process.}

\subsection{Connecting to Governance and Decision-Making}
\paragraph{The outputs of the risk assessment process should feed directly into governance and decision-making at every level of the organisation. At the operational level, risk assessment outputs inform day-to-day process and control decisions. At the management level, they inform resource allocation, policy development, and operational oversight. At the board level, they provide the information needed to fulfil governance responsibilities around risk oversight and strategic direction.}
\paragraph{For this connection to work effectively, the risk assessment process needs to produce outputs that are useful to decision-makers — not just comprehensive for their own sake. A risk register with hundreds of entries at a uniform level of detail may satisfy a documentation requirement but will not effectively inform board-level governance. The ability to synthesise risk information — identifying the most material risks and presenting them clearly and proportionately — is a genuinely important professional skill.}

\subsection{The Role of the Compliance and Legal Function}
\paragraph{Within this broader governance picture, the compliance and legal function typically plays a second-line role — as we explored in Chapter~\ref{chap:Risk Management Frameworks}. This means:}
\begin{itemize}
    \item{\textbf{Supporting and challenging} the first line's risk identification and assessment, rather than doing it on their behalf.}
    \item{\textbf{Maintaining oversight} of the risk assessment process — ensuring it is being conducted consistently, thoroughly, and in accordance with regulatory requirements.}
    \item{\textbf{Providing expertise} on the legal and regulatory dimensions of identified risks.}
    \item{\textbf{Escalating} material risks to senior management and the board where necessary.}
    \item{\textbf{Evidencing compliance} — ensuring that the documentation produced by the risk assessment process is adequate to demonstrate regulatory compliance.}
\end{itemize}
\paragraph{This is a position of significant influence and significant responsibility. Done well, it means that the compliance and legal function genuinely shapes the organisation's risk profile and its ability to operate sustainably. Done poorly — or reduced to a box-ticking exercise — it creates a false sense of security that can prove very costly indeed.}

%---------------------------- SECTION ----------------------------
\section{A Word on Iteration}
\paragraph{Before we move into the detail of each step, it is worth saying something explicit about the iterative nature of the risk assessment process.}
\paragraph{The five steps described above are presented in sequence — and in practice, that is broadly the order in which they are carried out. But risk assessment is not a linear process in the sense of being completed once and then finished. It is a \textbf{cycle} — one that returns to the beginning as circumstances change, as controls are tested, and as new information comes to light.}
\paragraph{Moreover, within any given assessment, the steps interact and inform each other in ways that are not always predictable at the outset. The process of evaluating risks in Step Three may reveal hazards that were not identified in Step One. The process of designing controls in Step Four may uncover exposures that were not apparent in Step Two. A review conducted under Step Five may prompt a fundamental reassessment of the scope and criteria established before the process began.}
\paragraph{This iterative, dynamic character is not a weakness of the risk assessment process — it is a strength. It reflects the reality that risk is complex, that knowledge is imperfect, and that the goal is not to produce a perfect document but to develop a progressively more accurate and useful understanding of the organisation's risk landscape.}
\paragraph{The compliance professional who understands this — and who approaches risk assessment as an ongoing, evolving practice rather than a periodic exercise — is considerably better placed to deliver genuine value from the process.}

%---------------------------- SECTION ----------------------------
\newpage
\section{Chapter Summary}
In this chapter we learned about the following process elements:
\begin{table}[H]
  \centering
  \renewcommand{\arraystretch}{1.5}
  \begin{tabularx}{\textwidth}{ | >{\raggedright\arraybackslash}p{0.35\textwidth} 
                                | >{\raggedright\arraybackslash}X | }
    \hline
    \textbf{Element} & \textbf{Purpose} \\
    \hline
    \textbf{Scope} & Defines what the assessment covers — and what it does not\\
    \hline
    \textbf{Context} & Ensures the assessment reflects the organisation's real environment and objectives\\
    \hline
    \textbf{Risk criteria} & Establishes the standards against which risks will be evaluated\\
    \hline
    \textbf{Communication and consultation} & Builds completeness, credibility, and ownership throughout the process\\
    \hline
    \textbf{Governance} & Ensures clear accountability and connection to organisational decision-making\\
    \hline
    \textbf{Step 1 — Identify} & What could go wrong?\\
    \hline
    \textbf{Step 2 — Determine exposure} & Who or what is affected, and how?\\
    \hline
    \textbf{Step 3 — Evaluate} & How likely and how serious?\\
    \hline
    \textbf{Step 4 — Control} & What are we going to do about it?\\
    \hline
    \textbf{Step 5 — Record, review, monitor} & How do we document, maintain, and evidence the process?\\
    \hline
  \end{tabularx}
  \caption{Process Elements}
\end{table}

\paragraph{Key takeaways:}
\begin{itemize}
  \item{The risk assessment process follows a logical five-step cycle — but it is iterative and dynamic, not linear and finite.}
  \item{The preparatory work — defining scope, understanding context, and establishing risk criteria — is as important as the steps themselves.}
  \item{Communication, consultation, and governance run alongside the entire process — not as discrete steps but as continuous activities.}
  \item{The process needs to connect to organisational strategy, governance, and decision-making to deliver genuine value.}
  \item{The compliance and legal function plays a second-line role — supporting, challenging, overseeing, and evidencing rather than owning the risks themselves.}
  \item{The goal is not a perfect document — it is a progressively more accurate and useful understanding of the organisation's risk landscape.}
\end{itemize}

\barquote{In Chapter~\ref{chap:Step 1 — Identifying Hazards and Risk Sources}, we move into the first step of the process in detail — exploring the techniques, tools, and practical considerations involved in identifying hazards and risk sources as comprehensively and effectively as possible.}